% Source: https://artofproblemsolving.com/community/c3h1382138_algebra_problem
% Stable identity: aops-topic-1382138-q01
Let $t=\sqrt x\ge0$. Then $t^2+t-C=0$, so
\[
t=\dfrac{-1\pm\sqrt{1+4C}}{2}.
\]
The function $t\mapsto t^2+t$ is increasing for $t\ge0$ and takes every value $C\ge0$ exactly once, and no negative value. So there is a solution exactly when $C\ge0$, and then we need the root with the plus sign (the other root is negative):
\[
\sqrt x=\dfrac{\sqrt{1+4C}-1}{2},\qquad
x=\left(\dfrac{\sqrt{1+4C}-1}{2}\right)^2=\boxed{\dfrac{2C+1-\sqrt{4C+1}}{2}}\qquad(C\ge0).
\]
For example, $C=6$ gives $x=\dfrac{13-5}{2}=4$, and indeed $4+\sqrt4=6$.
